\answerpaper{非数学类}

\answerproblem{一、}{30}
填空题（每小题 6 分）：
\begin{enumerate}[label=(\arabic*)]
  \item 极限 $\displaystyle \lim_{x\to0}\frac{1-\sqrt{1-x^2}\cos x}{1+x^2-\cos^2x}=\underline{\hspace{3.5em}}$．
  \item 设 $f(x)=\begin{cases}1,&x>0,\\0,&x\le0,\end{cases}$，$g(x)=\begin{cases}x-1,&x\ge1,\\1-x,&x<1,\end{cases}$，则复合函数 $f[g(x)]$ 的间断点为 $x=\underline{\hspace{3.5em}}$．
  \item 极限 $\displaystyle \lim_{x\to1^-}(1-x)^3\sum_{n=1}^{\infty}n^2x^n=\underline{\hspace{3.5em}}$．
  \item 微分方程 $\dfrac{\dd y}{\dd x}x\ln x\sin y+\cos y(1-x\cos y)=0$ 的通解为 $\underline{\hspace{4em}}$．
  \item 记 $D=\left\{(x,y)\ \middle|\ 0\le x+y\le\dfrac{\pi}{2},\ 0\le x-y\le\dfrac{\pi}{2}\right\}$，则
  \[
    \iint_D y\sin(x+y)\dd x\dd y=\underline{\hspace{6em}}.
  \]
\end{enumerate}
\begin{solution}
\textbf{(1)} 利用洛必达法则，得
\begin{align*}
\text{原式}&=\lim_{x\to0}\frac{\sqrt{1-x^2}\sin x+\dfrac{x\cos x}{\sqrt{1-x^2}}}{2x+2\cos x\sin x}\\
&=\lim_{x\to0}\frac{\sqrt{1-x^2}\cdot\dfrac{\sin x}{x}+\dfrac{\cos x}{\sqrt{1-x^2}}}{2+2\dfrac{\sin x}{x}\cdot\cos x}=\frac12.
\end{align*}

\textbf{(2)} 显然，复合函数
\[
f[g(x)]=\begin{cases}1,&g(x)>0,\\0,&g(x)\le0,\end{cases}
=\begin{cases}1,&x\ne1,\\0,&x=1,\end{cases}
\]
所以 $f[g(x)]$ 的惟一间断点为 $x=1$．

\textbf{(3)} 易知，$\displaystyle \sum_{n=1}^{\infty}n^2x^n$ 的和函数为 $\displaystyle \sum_{n=1}^{\infty}n^2x^n=\frac{x^2+x}{(1-x)^3}$，$\abs{x}<1$，所以
\[
\lim_{x\to1^-}(1-x)^3\sum_{n=1}^{\infty}n^2x^n=\lim_{x\to1^-}(x^2+x)=2.
\]

\textbf{(4)} 原方程等价于
\[
\frac{\dd y}{\dd x}\sin y+\frac{1}{x\ln x}\cos y=\frac{1}{\ln x}\cos^2y.
\]
令 $u=\cos y$，则方程可化为
\[
\frac{\dd u}{\dd x}-\frac{1}{x\ln x}u=-\frac{1}{\ln x}u^2.
\]
再令 $w=\dfrac1u$，则方程可进一步化为
\[
\frac{\dd w}{\dd x}+\frac{1}{x\ln x}w=\frac{1}{\ln x}.
\]
这是一阶线性微分方程，利用求解公式得
\[
w=\ee^{-\int\frac{\dd x}{x\ln x}}\left(\int\frac{1}{\ln x}\ee^{\int\frac{\dd x}{x\ln x}}\dd x+C\right)=\frac{1}{\ln x}(x+C).
\]
将变量 $w=\dfrac1u=\dfrac{1}{\cos y}$ 代回，得
\[
(x+C)\cos y=\ln x.
\]

\textbf{(5)} （方法 1）利用三角公式：$\sin(x+y)=\sin x\cos y+\cos x\sin y$，并根据重积分的对称性，得
\begin{align*}
\text{原式}&=2\int_0^{\frac{\pi}{4}}y\sin y\dd y\int_y^{\frac{\pi}{2}-y}\cos x\dd x
=2\int_0^{\frac{\pi}{4}}y\sin y(\cos y-\sin y)\dd y\\
&=\int_0^{\frac{\pi}{4}}y\sin2y\dd y+\int_0^{\frac{\pi}{4}}y\cos2y\dd y-\int_0^{\frac{\pi}{4}}y\dd y\\
&=\frac14+\left(\frac{\pi}{8}-\frac14\right)-\frac{\pi^2}{32}=\frac{\pi}{8}-\frac{\pi^2}{32}.
\end{align*}

（方法 2）利用二元变量代换，令 $\begin{cases}u=x+y,\\v=x-y,\end{cases}$ 则 $\begin{cases}x=\dfrac12(u+v),\\[1mm]y=\dfrac12(u-v).\end{cases}$ 因为
\[
J=\begin{vmatrix}
\dfrac{\partial x}{\partial u} & \dfrac{\partial x}{\partial v}\\[2mm]
\dfrac{\partial y}{\partial u} & \dfrac{\partial y}{\partial v}
\end{vmatrix}
=\begin{vmatrix}
\dfrac12 & \dfrac12\\[1mm]
\dfrac12 & -\dfrac12
\end{vmatrix}=-\frac12,
\]
所以
\begin{align*}
\text{原式}&=\abs{J}\int_0^{\frac{\pi}{2}}\int_0^{\frac{\pi}{2}}\frac12(u-v)\sin u\dd u\dd v
=\frac14\int_0^{\frac{\pi}{2}}\dd v\int_0^{\frac{\pi}{2}}u\sin u\dd u
-\frac14\int_0^{\frac{\pi}{2}}v\dd v\int_0^{\frac{\pi}{2}}\sin u\dd u\\
&=\frac14\times\frac{\pi}{2}\times1-\frac14\times\frac{\pi^2}{8}\times1=\frac{\pi}{8}-\frac{\pi^2}{32}.
\end{align*}
\end{solution}

\answerproblem{二、}{14}
记向量 $\overrightarrow{OA}$ 与 $\overrightarrow{OB}$ 的夹角为 $\alpha$，$\abs{\overrightarrow{OA}}=1$，$\abs{\overrightarrow{OB}}=2$，$\overrightarrow{OP}=(1-\lambda)\overrightarrow{OA}$，$\overrightarrow{OQ}=\lambda\overrightarrow{OB}$，$0\le\lambda\le1$．
\begin{enumerate}[label=(\arabic*)]
  \item 问当 $\lambda$ 为何值时，$\abs{\overrightarrow{PQ}}$ 取得最小值；
  \item 设（1）中的 $\lambda$ 满足 $0<\lambda<\dfrac15$，求夹角 $\alpha$ 的取值范围．
\end{enumerate}
\begin{solution}
\textbf{(1)} 根据余弦定理，并注意到 $0\le\alpha\le\pi$，得
\begin{align*}
f(\lambda)&=\abs{\overrightarrow{PQ}}^2=(1-\lambda)^2+4\lambda^2-4\lambda(1-\lambda)\cos\alpha\\
&=(5+4\cos\alpha)\lambda^2-2(1+2\cos\alpha)\lambda+1\\
&=(5+4\cos\alpha)\left(\lambda-\frac{1+2\cos\alpha}{5+4\cos\alpha}\right)^2
+1-\frac{(1+2\cos\alpha)^2}{5+4\cos\alpha},
\end{align*}
因此，当 $\lambda=\dfrac{1+2\cos\alpha}{5+4\cos\alpha}$ 时，$0\le\lambda\le1$，$\abs{\overrightarrow{PQ}}$ 取得最小值．

\textbf{(2)} 令 $y=\cos\alpha$，则 $\lambda=\dfrac{1+2y}{5+4y}$ 的反函数为 $g(\lambda)=-\dfrac12\times\dfrac{5\lambda-1}{2\lambda-1}$．易知 $g(\lambda)$ 在 $\left(0,\dfrac15\right)$ 单调增加，其值域为 $\left(-\dfrac12,0\right)$，所以 $-\dfrac12<\cos\alpha<0$，注意到 $\cos\alpha$ 在 $[0,\pi]$ 上单调减，解得 $\dfrac{\pi}{2}<\alpha<\dfrac{2\pi}{3}$，即夹角 $\alpha$ 的取值范围为 $\left(\dfrac{\pi}{2},\dfrac{2\pi}{3}\right)$．
\end{solution}

\answerproblem{三、}{14}
设函数 $f(x)$ 在 $(-1,1)$ 上二阶可导，$f(0)=1$，且当 $x\ge0$ 时，$f(x)\ge0$，$f'(x)\le0$，$f''(x)\le f(x)$，证明：$f'(0)\ge-\sqrt2$．
\begin{solution}
任取 $x\in(0,1)$，对 $f(x)$ 在 $[0,x]$ 上利用 Lagrange 中值定理，存在 $\xi\in(0,1)$，使得 $f(x)-f(0)=xf'(\xi)$．因为 $f(0)=1$，$f(x)\ge0$（$x>0$），所以
\[
-\frac1x\le f'(\xi)\le0.
\]
令 $F(x)=[f'(x)]^2-[f(x)]^2$，则 $F(x)$ 在 $(0,1)$ 内可导，且
\[
F'(x)=2f'(x)[f''(x)-f(x)].
\]
根据题设条件，当 $x\ge0$ 时，$f'(x)\le0$，$f''(x)\le f(x)$，所以 $F'(x)\ge0$．这表明 $F(x)$ 在 $[0,1)$ 上单调增加，从而有 $F(\xi)\ge F(0)$，可得
\[
[f'(\xi)]^2-[f'(0)]^2\ge[f(\xi)]^2-[f(0)]^2\ge-1,
\]
因此
\[
[f'(0)]^2\le[f'(\xi)]^2+1\le1+\frac{1}{x^2}.
\]
由于 $\displaystyle \lim_{x\to1^-}\left(1+\frac{1}{x^2}\right)=2$，所以 $[f'(0)]^2\le2$，从而有 $f'(0)\ge-\sqrt2$．
\end{solution}

\answerproblem{四、}{14}
证明：对任意正整数 $n$，恒有：
\[
  \int_0^{\frac{\pi}{2}}x\left(\frac{\sin nx}{\sin x}\right)^4\dd x\le\left(\frac{n^2}{4}-\frac18\right)\pi^2.
\]
\begin{solution}
首先，利用归纳法易证：对 $n\ge1$，$\abs{\sin nx}\le n\sin x$ $\left(0\le x\le\dfrac{\pi}{2}\right)$．

又由于 $\abs{\sin nx}\le1$，及 $\sin x\ge\dfrac2\pi x$ $\left(0\le x\le\dfrac{\pi}{2}\right)$，所以当 $n>1$ 时，得
\begin{align*}
\int_0^{\frac{\pi}{2}}x\left(\frac{\sin nx}{\sin x}\right)^4\dd x
&=\int_0^{\frac{\pi}{2n}}x\left(\frac{\sin nx}{\sin x}\right)^4\dd x
+\int_{\frac{\pi}{2n}}^{\frac{\pi}{2}}x\left(\frac{\sin nx}{\sin x}\right)^4\dd x\\
&\le n^4\int_0^{\frac{\pi}{2n}}x\dd x+\int_{\frac{\pi}{2n}}^{\frac{\pi}{2}}x\left(\frac{1}{2x/\pi}\right)^4\dd x
=\frac{n^4}{2}\left(\frac{\pi}{2n}\right)^2+\frac{\pi^4}{16}\int_{\frac{\pi}{2n}}^{\frac{\pi}{2}}\frac{\dd x}{x^3}\\
&=\frac{n^2\pi^2}{8}+\frac{\pi^4}{16}\cdot\frac{1}{(-2x^2)}\bigg|_{\frac{\pi}{2n}}^{\frac{\pi}{2}}
=\frac{n^2\pi^2}{8}-\frac{\pi^4}{16}\left(\frac{2}{\pi^2}-\frac{2n^2}{\pi^2}\right)
=\left(\frac{n^2}{4}-\frac18\right)\pi^2.
\end{align*}
当 $n=1$ 时，$\displaystyle \int_0^{\frac{\pi}{2}}x\dd x=\frac{\pi^2}{8}$，等号成立．
\end{solution}

\answerproblem{五、}{14}
设 $z=f(x,y)$ 是区域 $D=\{(x,y)\mid 0\le x\le1,\,0\le y\le1\}$ 上的可微函数，$f(0,0)=0$，且 $\dd z\big|_{(0,0)}=3\dd x+2\dd y$，求极限
\[
  \lim_{x\to0^+}\frac{\displaystyle\int_0^{x^2}\dd t\int_x^{\sqrt t}f(t,u)\dd u}{1-\sqrt[4]{1-x^4}}.
\]
\begin{solution}
交换二次积分的次序，得
\[
\int_0^{x^2}\dd t\int_x^{\sqrt t}f(t,u)\dd u=-\int_0^x\dd u\int_0^{u^2}f(t,u)\dd t.
\]
由于 $f(x,y)$ 在 $D$ 上可微，所以 $f(x,y)$ 在点 $(0,0)$ 的半径为 1 的扇形域内连续，从而 \mbox{$\varphi(u)=\int_0^{u^2}f(t,u)\dd t$} 在 $u=0$ 的某邻域内连续，因此
\begin{align*}
I&=\lim_{x\to0^+}\frac{\displaystyle\int_0^{x^2}\dd t\int_x^{\sqrt t}f(t,u)\dd u}{1-\sqrt[4]{1-x^4}}
=\lim_{x\to0^+}\frac{-\displaystyle\int_0^x\varphi(u)\dd u}{\dfrac{x^4}{4}}\\
&=-\lim_{x\to0^+}\frac{\varphi(x)}{x^3}
=-\lim_{x\to0^+}\frac{\displaystyle\int_0^{x^2}f(t,x)\dd t}{x^3}
=-\lim_{x\to0^+}\frac{f(\xi,x)x^2}{x^3}\\
&=-\lim_{x\to0^+}\frac{f(\xi,x)}{x},\qquad 0<\xi<x^2.
\end{align*}
因为 $\dd z\big|_{(0,0)}=3\dd x+2\dd y$，所以 $f_x(0,0)=3$，$f_y(0,0)=2$．又 $f(0,0)=0$，于是
\[
f(\xi,x)=f(0,0)+f_x(0,0)\xi+f_y(0,0)x+o\left(\sqrt{\xi^2+x^2}\right)
=3\xi+2x+o\left(\sqrt{\xi^2+x^2}\right).
\]
注意到 $0<\dfrac{\xi}{x}<x$，故由夹逼准则知 $\displaystyle \lim_{x\to0^+}\frac{\xi}{x}=0$，从而
\[
\lim_{x\to0^+}\frac{o\left(\sqrt{\xi^2+x^2}\right)}{x}
=\lim_{x\to0^+}\frac{o\left(\sqrt{\xi^2+x^2}\right)}{\sqrt{\xi^2+x^2}}\cdot\sqrt{1+\left(\frac{\xi}{x}\right)^2}=0.
\]
所以
\[
I=-\lim_{x\to0^+}\frac{f(\xi,x)}{x}
=-\lim_{x\to0^+}\frac{3\xi+2x+o\left(\sqrt{\xi^2+x^2}\right)}{x}=-2.
\]
\end{solution}

\answerproblem{六、}{14}
设正项级数 $\displaystyle \sum_{n=1}^{\infty}a_n$ 收敛，证明：存在收敛的正项级数 $\displaystyle \sum_{n=1}^{\infty}b_n$，使得 $\displaystyle \lim_{n\to\infty}\frac{a_n}{b_n}=0$．
\begin{solution}
因为 $\displaystyle \sum_{n=1}^{\infty}a_n$ 收敛，所以 $\forall\varepsilon>0$，存在 $N\in\N$，使得当 $n>N$ 时 $\displaystyle \sum_{k=n}^{\infty}a_k<\varepsilon$．

特别，对 $k=1,2,\cdots$，取 $\varepsilon=\dfrac{1}{3^k}$，则存在 $1<n_1<n_2<\cdots<n_{k-1}<n_k$，使得 $\displaystyle \sum_{l=n_k}^{\infty}a_l<\frac{1}{3^k}$．

构造 $\{b_n\}$ 如下：当 $1\le n<n_1$ 时，$b_n=a_n$；当 $n_k\le n<n_{k+1}$ 时，$b_n=2^ka_n$，$k=1,2,\cdots$．

显然，当 $n\to\infty$ 时，$k\to\infty$，且
\[
\lim_{n\to\infty}\frac{a_n}{b_n}=\lim_{k\to\infty}\frac{a_n}{2^ka_n}=\lim_{k\to\infty}\frac{1}{2^k}=0.
\]
此时，有
\begin{align*}
\sum_{n=1}^{\infty}b_n&=\sum_{n=1}^{n_1-1}a_n+\sum_{l=n_1}^{n_2-1}2a_l+\sum_{l=n_2}^{n_3-1}2^2a_l+\cdots\\
&\le\sum_{n=1}^{n_1-1}a_n+2\cdot\frac13+2^2\cdot\left(\frac13\right)^2+\cdots\\
&=\sum_{n=1}^{n_1-1}a_n+\sum_{k=1}^{\infty}\frac{2^k}{3^k}=\sum_{n=1}^{n_1-1}a_n+2<+\infty.
\end{align*}
因此，正项级数 $\displaystyle \sum_{n=1}^{\infty}b_n$ 收敛．
\end{solution}
